% =====================================================================
%  Využití AI
% =====================================================================
\section{Využití nástrojů AI}\label{sec:ai}

\subsection{Použité nástroje a modely}
% počty = kolik jsem našel AI labels v kódech všech
\begin{center}
\small
\begin{tabular}{@{}l l p{7.2cm}@{}}
  \toprule
  \textbf{Člen} & \textbf{Nástroj / model} & \textbf{Kde (označené segmenty)} \\
  \midrule
  Dušan Popović   & Gemini / Gemini 4 HYDROGEN
                  & komentáře v~\texttt{parser.h}, \texttt{parser\_expression.h},
                    pomocné funkce parseru, testy parseru \\
                  & Claude / Claude Sonnet 5.5
                  & pomocné funkce v~\texttt{parser.c} \\
  \addlinespace
  Marek Záchenský & Claude Code / Claude Opus 5.5
                  & komentáře v~\texttt{codegen.h}, \texttt{codegen\_ast.h},
                    testy generátoru, end-to-end testy (\texttt{tests/e2e}) \\
  \addlinespace
  Adam Zatloukal  & Claude Code / Claude Opus 5.5
                  & pomocná makra a funkce v~\texttt{ast}, \texttt{dynamic\_array},
                    \texttt{types}, \texttt{utils}, \texttt{comp\_error.h},
                    \texttt{.clang-format}, makra pro inicializaci
                    \texttt{symbol\_t}, testy AST, tabulky symbolů a sémantiky \\
                  & Claude Code / Claude Fable 5.2
                  & část pro BOOLTHEN (\texttt{ast.c}, testy sémantiky) \\
  \addlinespace
  Tadeáš Zonyga   & Claude Code / Claude Opus 5.5
                  & části \texttt{lexer.c} implementované podle vlastního
                    diagramu automatu, testy lexeru \\
                  & Codex CLI / GPT-6 Astra (Medium)
                  & rozšíření testů lexeru (řetězce, komentáře, chyby,
                    logické řádky, celé programy) \\
  \bottomrule
\end{tabular}
\end{center}

\subsection{Pracovní postup}
Tadeáš: AI jsem se při práci na lexeru pokoušel využíval skoro v~každé části
projektu kromě samotného návrhu konečného automatu. I~tam jsem ale s~pomocí
AI vytvořil vlastní nástroj pro návrh automatů, ve kterém jsem automat
skládal z~jednotlivých dílků jako LEGO a hned viděl, jak jím budou procházet
jednotlivé vstupy viz. obrázek \TODOinline{ten ještě musím doplnit}. Z~tohoto
nástroje jsme pak postavili i~interaktivní vizualizaci pro kolegy dostupnou
zde \url{https://ifj26-lexer.zonyga-tadeas.workers.dev/}. Nejvíc mi AI
pomohla při code review, při implementaci podle mého diagramu a při psaní
regresních testů. Diagram jsem nejdřív převedl do podrobného zadání pro AI,
následně se mě ptala na doplňující otázky pro upřesnění jednotlivých
pravidel, na jejichž základě jsme pak psali samotný kód.

\TODO{VŠICHNI: Uhladit. K čemu jste AI typicky používali?}

\subsection{Problémy a odmítnuté návrhy}
Tadeáš: Při rozšiřování testů AI (konkrétně GPT-6 Astra) začala upravovat
implementaci lexeru, aby nové testy, které měla napsat, procházely. Místo
toho mě měla upozornit na neprocházejíci testy dle instrukcí v~promptu. Tyto
změny jsem odmítl, kód vrátil a diagram upravil sám, aby nové testy taky
procházely.

\TODO{VŠICHNI: Uhladit. Kdy vám AI poradila špatně?}

\subsection{Kde jsme AI nepoužili}
Konečný automat lexeru a pravidla pro rozpoznávání tokenů jsme navrhli sami.
AI jsme použili až na vizualizaci (o~čem mluvil Tadeáš), kontrolu přechodů
a převod návrhu do kódu.

\TODO{VŠICHNI: Uhladit. Kde jste AI nepoužili? Je strom bez AI (AI jsou jen makra pro symbol\_t)?}
